% TOP_PROBLEM: {"id": 30007028, "title": "Generalized lc flip termination on Kähler fourfolds", "category_id": 5, "category": {"id": 5, "name": "algebraic_geometry", "display_name": "Algebraic Geometry", "description": "Geometric objects defined by polynomial equations.", "slug": "algebraic-geometry", "order_index": 5, "created_at": "2026-07-31T15:26:25.671Z"}, "status": "open"}
% FIRST_POSTED: 2026-10-06
% ENTRY_KIND: statement_only
% !TeX program = lualatex
\documentclass[11pt]{article}
\usepackage[margin=25mm]{geometry}
\usepackage{fontspec}
\setmainfont{Latin Modern Roman}
\usepackage{amsmath,amssymb,mathtools}
\usepackage{unicode-math}
\setmathfont{Latin Modern Math}
\usepackage{xurl,hyperref}
\hypersetup{colorlinks=true,urlcolor=blue}
\setcounter{secnumdepth}{0}
\setlength{\parindent}{0pt}
\setlength{\parskip}{5pt}
\setlength{\emergencystretch}{3em}
\newcommand{\R}{\mathbb R}
\newcommand{\N}{\mathbb N}
\newcommand{\Z}{\mathbb Z}
\newcommand{\Q}{\mathbb Q}
\newcommand{\C}{\mathbb C}
\newcommand{\D}{\mathbb D}
\newcommand{\F}{\mathbb F}
\newcommand{\sref}[2]{\hyperlink{source:#1:#2}{[S#2]}}
\newcommand{\cataloguescope}[1]{#1}
\begin{document}
\section*{Generalized lc flip termination on Kähler fourfolds}\label{30007028}
\subsection{Definitions and mathematical statement}
\paragraph{Statement recorded in the supplied source.}
On globally Weil-Q-factorial normal compact Kähler fourfolds, fix effective rational boundary and a nef rational b-Cartier divisor represented on a projective birational Kähler model, with the generalized pair lc. Must every sequence of existing projective small connected-fiber flips for this same rational adjoint and b-data terminate? Each flipping adjoint is negative ample, and its next transform positive ample.
\subsection{Short English statement}
On globally Weil-Q-factorial normal compact Kähler fourfolds, fix effective rational boundary and a nef rational b-Cartier divisor represented on a projective birational Kähler model, with the generalized pair lc. Must every sequence of existing projective small connected-fiber flips for this same rational adjoint and b-data terminate? Each flipping adjoint is negative ample, and its next transform positive ample.
\subsection{Sources}
\begin{itemize}
\item[S1] Original question source or mathematical context.
\url{http://aimpl.org/minimalmodcharp/1/}.
\item[S2] Original question source or mathematical context.
\url{https://arxiv.org/abs/2209.13122}.
\item[S3] Original question source or mathematical context.
\url{https://epiga.episciences.org/16998}.
\item[S4] Original question source or mathematical context.
\url{https://jep.centre-mersenne.org/item/10.5802/jep.334.pdf}.
\item[S5] Original question source or mathematical context.
\url{https://master-math-fonda.imj-prg.fr/gt/mmp.pdf}.
\item[S6] OpenAI, manuscript citation page. Source of the formulation motivating this local problem.
\url{https://github.com/openai/math/blob/adc7f1241b42e322a6451854ab7e4b4c146bf78a/preprints/Termination-of-generalized-log-canonical-flips-on-compact-Kahler-fourfolds-October-5-2026/README.md}.
\end{itemize}
\end{document}
